% ==================================================================
% id:       prob-010
% topic:    מבוא להסתברות
% subtopic: תכונות של פונקציית הסתברות
% solution: full
% ==================================================================
באוכלוסייה מסוימת 40\% הן נשים ו-50\% הם בעלי רישיון נהיגה. ידוע ש-40\% מהאוכלוסייה הם גברים בעלי רישיון נהיגה. בוחרים אדם באקראי מהאוכלוסייה.
\begin{enumerate}
\item מה ההסתברות שזו אישה בעלת רישיון נהיגה?
\begin{solbox}
נסמן $W$ = „אישה”, $L$ = „בעל/ת רישיון”. $L$ הוא איחוד זר של $W\cap L$ ושל $W^C\cap L$, ולכן $\PP(W\cap L)=\PP(L)-\PP(W^C\cap L)=0.5-0.4=\ans{0.1}$.
\end{solbox}
\item מה ההסתברות שזה גבר שאין לו רישיון נהיגה?
\begin{solbox}
$\PP(W^C)=0.6$, ו-$W^C$ הוא איחוד זר של $W^C\cap L$ ושל $W^C\cap L^C$. לכן $\PP(W^C\cap L^C)=0.6-0.4=\ans{0.2}$.
\end{solbox}
\end{enumerate}
